% BHU PYQ -- Professional Exam Template
\documentclass[11pt,a4paper]{article}

\usepackage[a4paper,top=2.5cm,bottom=2.5cm,left=2.8cm,right=2.8cm,headheight=14pt]{geometry}
\usepackage{amsmath,amssymb}
\usepackage[T1]{fontenc}
\usepackage[utf8]{inputenc}
\usepackage{lmodern}
\usepackage{microtype}
\usepackage{fancyhdr}
\usepackage{enumitem}
\usepackage{listings}
\usepackage{hyperref}
\usepackage{lastpage}
\usepackage{parskip}

\hypersetup{colorlinks=true,urlcolor=black,linkcolor=black,
  pdftitle={BCS-504A: System Analysis Design -- BHU PYQ}}

\pagestyle{fancy}
\fancyhf{}
\renewcommand{\headrulewidth}{0.4pt}
\renewcommand{\footrulewidth}{0.4pt}
\fancyhead[L]{\small   \textit{Banaras Hindu University}}
\fancyhead[R]{\small   \textit{BCS-504A: System Analysis Design}}
\fancyfoot[C]{\small Page  \thepage\ of \pageref{LastPage}}

\lstset{
  basicstyle=\small tfamily,
  frame=single,
  breaklines=true,
  columns=flexible,
  keepspaces=true
}


\newlist{parts}{enumerate}{2}
\setlist[parts,1]{label=(\alph*),leftmargin=*,itemsep=4pt,topsep=3pt}
\setlist[parts,2]{label=(\roman*),leftmargin=*,itemsep=2pt,topsep=2pt}

\newcommand{\pts}[1]{\hfill   \textit{[#1]}}


\begin{document}

\begin{center}
  {\large\bfseries BANARAS HINDU UNIVERSITY}\\[2pt]
  {
\normalsize B.Sc. (Hons.) Semester V Examination, 2016-17}\\[6pt]
  \rule{0.8\linewidth}{0.5pt}\\[6pt]
  {\large\bfseries Paper: BCS-504A --- System Analysis Design}\\[4pt]
  {
\normalsize Time: Three hours \hfill Maximum Marks: 70}
\end{center}

\rule{\linewidth}{0.4pt}

\smallskip



\noindent   \textit{Instructions: ANSWER ANY FIVE QUESTIONS AND THE FIGURES IN THE RIGHT HAND MARGIN INDICATE:}

\bigskip

\medskip


\noindent   \textbf{Question 1.}

\begin{parts}
  \item Write the IEEE definitions of the following terms in connection with software: \pts{6}
  
\begin{parts}
    \item Error
    \item Fault
    \item Failure
  \end{parts}
  \item Design an ER diagram for a Library Management System. Specify all constraints that should hold on the database. Make sure that it has at least three entity types, two relationship types, a weak entity type, a superclass/subclass relationship, and an $n$-ary ($n > 2$) relationship type. \pts{8}
\end{parts}

\medskip


\noindent   \textbf{Question 2.}

\begin{parts}
  \item Describe software crisis and the goals of software engineering. \pts{7}
  \item Write four important attributes of software products and differentiate them from hardware products. \pts{7}
\end{parts}

\medskip


\noindent   \textbf{Question 3.}

\begin{parts}
  \item Compare prototyping and iterative software development process models by listing their strengths and weaknesses. \pts{7}
  \item Describe at least five important characteristics of an SRS. \pts{7}
\end{parts}

\medskip


\noindent   \textbf{Question 4.}

\begin{parts}
  \item What is the role of problem analysis in software development? Explain the use of Data Flow Diagrams in problem analysis with an example. \pts{7}
  \item What are use cases? Where are these use cases used during the software development process? Explain through an example. \pts{7}
\end{parts}

\medskip


\noindent   \textbf{Question 5.}

\begin{parts}
  \item What do you understand by a module? What are the modularization criteria used for function-oriented design to produce modular designs? \pts{7}
  \item Explain all four major steps of structured design methodology using one example. \pts{7}
\end{parts}

\medskip


\noindent   \textbf{Question 6.}

\begin{parts}
  \item Write the purpose of the following diagrams of UML used in object-oriented design: \pts{6}
  
\begin{parts}
    \item Class diagram
    \item Interaction diagram
    \item Activity diagram
  \end{parts}
  \item What do you understand by test cases and test criteria? Write the differences between black-box and white-box testing. \pts{8}
\end{parts}

\medskip


\noindent   \textbf{Question 7.} Write short notes on the following: \pts{14}

\begin{parts}
  \item Test oracle
  \item Software maintenance
  \item Data dictionary
  \item Open-closed principle
\end{parts}

\vfill
\rule{\linewidth}{0.4pt}

\begin{center}\small
\href{https://bhu-exam.vercel.app/}{Click here to visit our website}\\[4pt]\href{https://me-aryan.vercel.app/}{Click here to visit our contributor}\\[4pt]\href{https://quashan-me.vercel.app/}{Click here to visit our contributor}
\end{center}
\end{document}